%% ============================================================
%% 分片 31：原文第 11 节「$A$-膜与 ${\cal D}$-模」（原文 glshort.tex 起于第 9938 行）
%% 由 tools/harvmac2tex.py 从 arXiv 源码（hep-th/0604151v3）自动生成。
%% 只译英文散文；公式、编号、\label、\cite 一律不改。
%% ============================================================

\section{$A$-膜与 ${\cal D}$-模}
\label{sec:abranes}

\begin{keytakeaways}
  \begin{itemize}[leftmargin=1.4em,itemsep=2pt,topsep=2pt]
    \item \textbf{物理图像}：几何 Langlands 对应的左端已知是电特征膜（即 $\pi_1(C)\to{}^L\negthinspace G$ 的同态），其对偶是缠在 Hitchin 纤维上、带平坦 Chan-Paton 丛的磁特征膜；但通常表述的右端并不是 $\MH$ 上的 $(A,B,A)$ 型膜，而是 $G$-丛模空间上的 ${\cal D}$-模。本节的物理任务就是补上这最后一段翻译。
    \item \textbf{数学工具}：${\cal D}$-模是微分算子层的模；把 $\M(G,C)$ 上的 ${\cal D}$-模看成其余切丛的非交换形变上的 $B$-膜，则 $A$-膜的范畴与这类 $B$-膜范畴的等价（原文引 Kapustin 的论证）使「$A$-膜给出 ${\cal D}$-模」成为该等价的一个特例。
    \item \textbf{关键结论}：$\MH(G,C)$ 上复结构 $K$ 中的每个 $A$-膜都自动给出（粗略地说）$\M(G,C)$ 上的一个 ${\cal D}$-模；这就是所谓 A-膜到 D-模的翻译器，是几何 Langlands 纲领与规范理论对偶之间缺失的一环。
    \item \textbf{边界与展望}：严格处理不能只停留在稳定丛的模空间 $\M(G,C)$，而要在所有 $G$-丛的叠 ${\rm Bun}_G(C)$ 上构造 ${\cal D}$-模（Hecke 算子的作用方式要求如此）；相应地，所关心的膜必须能在底层规范理论中定义，而不只是在靶空间为 $\MH(G,C)$ 的 sigma 模型中定义，这一点在第~\ref{sec:zapendix}~节证明。
  \end{itemize}
\end{keytakeaways}


迄今所讲的内容还缺一个关键要素。按照几何 Langlands 纲领，$\pi_1(C)$ 到 $^L\negthinspace G$ 的同态与 ${\EUM}(G,C)$ 上的 Hecke 特征层之间存在自然的对应，而这些特征层同时也是 ${\cal D}$-模。所谓 ${\cal D}$-模，就是微分算子层的模；这个概念下面还会进一步说明。

至于对应式的左端，我们在第~\ref{sec:electricbranes}~节已经知道：$\pi_1(C)$ 到 $^L\negthinspace G$ 的同态对应于电特征膜。它们的对偶因此必定是磁特征膜；这些对偶是 ${\CMF}$ 型膜，即缠在 Hitchin 纤维化的一条纤维上、带平坦 Chan-Paton 丛的膜。这一陈述可由 $S$-对偶得到，也可以借助谱覆盖构造所提供的阿贝尔化直接验证 \cite{ref26}。

然而，这样得到的磁特征膜是 $\MH(G,C)$ 上的 $(A,B,A)$ 型膜；而在几何 Langlands 对应的通常表述中，右端应当涉及的却是（粗略地说）$G$-丛的模空间 $\M(G,C)$ 上的 ${\cal D}$-模。

本节的任务就是补上这缺失的一环。我们将证明：$\MH(G,C)$ 上复结构 $K$ 中的每一个 $A$-膜都自动给出（粗略地说）$\M(G,C)$ 上的一个 ${\cal D}$-模。

人们可以把 $\M(G,C)$ 上的 ${\cal D}$-模看作 $\M(G,C)$ 的余切丛的非交换形变上的 $B$-膜。此前已有论证 \cite{ref41}：对某一类全纯辛流形，$A$-膜的范畴等价于同一流形的非交换形变上 $B$-膜的范畴。$\MH(G,C)$ 上的 $A$-膜与 $\M(G,C)$ 上的 ${\cal D}$-模之间的联系，可以看作这一点的一个特例。

事实上，在几何 Langlands 纲领中，只在稳定丛的模空间 $\M(G,C)$ 上工作是不够的（而且由于 Hecke 算子的作用方式，也不可能这样做）。人们改而在 $C$ 上所有 $G$-丛构成的「叠」 ${\rm Bun}_G(C)$ 上构造 ${\cal D}$-模。所有 $G$-丛构成的叠在规范理论中的类比物，就是所有联络构成的空间。通常的说法是「${\cal D}$-模定义在 ${\rm Bun}_G(C)$ 上」；在我们的构造中与之对应的说法是：所关心的膜可以在底层的规范理论中定义，而不只是在靶空间为 $\MH(G,C)$ 的 sigma 模型中定义。对重要的膜证明这一点，实际上正是第~\ref{sec:zapendix}~节的目标。

\begin{yizhu}[含糊点拨]
按：本片两处「粗略地说」（roughly speaking）值得留意。一处修饰「$\M(G,C)$ 上的 ${\cal D}$-模」：严格的对象其实是 ${\rm Bun}_G(C)$ 上的（扭化）${\cal D}$-模，只因 Hecke 算子会把稳定丛变出稳定丛之外，才不能只停在 $\M(G,C)$ 上。另一处修饰「每个 $A$-膜都自动给出一个 ${\cal D}$-模」：翻译器的精确定义还要选取合适的 $A$-膜范畴（含余各向同性的对象）并对 ${\cal D}$-模作相应的扭化，所以原文才说「粗略地说」。
\end{yizhu}

\begin{yizhu}[易错点]
按：本片出现三个容易混淆的空间：$\MH(G,C)$ 是超凯勒的 Hitchin 模空间，膜（包括磁特征膜）住在这里，其复结构 $K$ 中的 $A$-膜是本片的主角；$\M(G,C)$ 是稳定 $G$-丛的模空间，${\cal D}$-模粗略地说住在这里；${\rm Bun}_G(C)$ 则是所有 $G$-丛构成的叠，才是 ${\cal D}$-模严格的家。镜像对称的三角中三条边分别是：$\MH({}^L\negthinspace G)$ 上的 $A$-膜、$\MH(G,C)$ 上的 $B$-膜、以及 ${\rm Bun}_G(C)$ 上的 ${\cal D}$-模；本片补的正是「$A$-膜 $\to$ ${\cal D}$-模」这一条边。
\end{yizhu}

\begin{yizhu}[物理图像]
按：把 ${\cal D}$-模想成余切丛上的 $B$-膜是理解本片的关键：$\M$ 的局部坐标 $x^i$ 与其对偶 $p_i$ 若实现为 $x^i$ 与 $\partial/\partial x^i$，就不再对易，于是 $T^*\M$ 变成一个非交换空间，而 ${\cal D}$-模恰是它上面的「$B$-膜」。$\MH$ 包含 $T^*\M^{st}$ 作为 Zariski 开集，因此 $\MH$ 上的 $A$-膜与 $\M$ 上的 ${\cal D}$-模之间的关系自然被期待为「同一块非交换几何的两个侧面」；本片正是把这一期待落实为具体构造的起点。
\end{yizhu}

\begin{yizhu}[物理图像]
按：「所有 $G$-丛构成的叠」在规范理论中的类比物是「所有联络构成的空间」，这句话解释了对偶为什么必须离开半稳定丛：'t Hooft 算子（Hecke 算子）在四维规范理论中作用于所有联络，而不只作用于满足稳定性条件的解；因此「${\cal D}$-模定义在 ${\rm Bun}_G(C)$ 上」这一数学事实，对应的是「所关心的膜可以在底层规范理论中定义，而不只是在靶空间为 $\MH(G,C)$ 的 sigma 模型中定义」。
\end{yizhu}

\begin{yizhu}[含糊点拨]
按：原文说 $A$-膜范畴与「某类全纯辛流形的非交换形变上 $B$-膜范畴」的等价「此前已有论证」，这是指 Kapustin 关于 A-膜与非交换几何的工作（本片此处的引用即该文）。注意原文用的动词是「argued（论证）」而非「proved（证明）」：这一等价在原文写作时还只是物理论证，并非定理；若读者要引用严格结果，须另找关于这类范畴等价的数学文献。
\end{yizhu}
